% Documento de trabajo: Esquema restructurado de los 8 temas
% Álgebra lineal en computación cuántica - UNIR
% -------------------------------------------------------
% Este fichero no es un tema completo: es un documento de planificación
% que recoge el esquema (esquemaExplorador) propuesto para cada tema,
% junto con la motivación cuántica que encabeza cada uno.
% -------------------------------------------------------

\documentclass[]{unirTema}
\input{personalizacion}

\begin{document}

\author{Francisco Costa Cano}
\color{negrounir}
\titulacion{Máster en computación cuántica}
\asignatura{Álgebra lineal en computación cuántica}

% ============================================================
%  BLOQUE 1: FUNDAMENTOS MATEMÁTICOS
% ============================================================
\bloque{1}{Fundamentos matemáticos}

% ------------------------------------------------------------
\tema{1}{Números complejos y su geometría}
% ------------------------------------------------------------

\portada

\begin{info}
  \textbf{Motivación cuántica.}
  Las amplitudes de probabilidad que gobiernan los sistemas cuánticos son números complejos.
  Antes de enunciar los postulados de la mecánica cuántica de forma precisa, necesitamos
  dominar el lenguaje matemático en que están escritos.
  Este tema proporciona esa base: los números complejos y su geometría.
\end{info}

\begin{esquemaExplorador}
  \temaEsquema{Forma binómica}{
    \conceptoEsquema{Definición de número complejo}{$z = a + bi,\quad a,b\in\R$}
    \conceptoEsquema{Conjugado complejo}{$\bar{z} = a - bi$}
    \conceptoEsquema{Módulo}{$|z| = \sqrt{a^2+b^2}$}
    \conceptoEsquema{Operaciones: suma, producto y cociente}{}
  }
  \temaEsquema{Geometría del plano complejo}{
  \conceptoEsquema{Representación como punto o vector en $\R^2$}{}
  \conceptoEsquema{Forma polar}{$z = |z|(\cos\theta + i\sin\theta)$}
  \conceptoEsquema{Forma exponencial (fórmula de Euler)}{$z = |z|e^{i\theta}$}
  \conceptoEsquema{Multiplicación como rotación y escalado}{}
  }
  \temaEsquema{Fórmula de De Moivre y raíces}{
    \conceptoEsquema{Potencias}{$z^n = |z|^n e^{in\theta}$}
    \conceptoEsquema{Raíces $n$-ésimas de la unidad}{$\omega_k = e^{2\pi ik/n}$}
  }
  \temaEsquema{Amplitud cuántica}{
    \conceptoEsquema{Magnitud y fase de una amplitud}{}
    \conceptoEsquema{Interferencia constructiva y destructiva}{}
    \conceptoEsquema{Interpretación probabilística: $|\alpha|^2$}{}
  }
\end{esquemaExplorador}

% ------------------------------------------------------------
\tema{2}{Espacios vectoriales complejos}
% ------------------------------------------------------------

\portada

\begin{info}
  \textbf{Motivación cuántica (anticipo del Postulado~I).}
  El Postulado~I de la mecánica cuántica establece que el estado de un sistema cuántico
  queda completamente descrito por un vector en un espacio vectorial complejo.
  Este tema construye ese espacio desde sus axiomas, proporcionando el marco en que
  vivirán los estados cuánticos a lo largo de todo el máster.
\end{info}

\begin{esquemaExplorador}
  \temaEsquema{Postulado I: anticipo}{
    \conceptoEsquema{El estado cuántico como vector: $\ket{\psi}\in\mathcal{H}$}{}
    \conceptoEsquema{Superposición cuántica}{$\alpha\ket{0}+\beta\ket{1}$}
    \conceptoEsquema{¿Qué estructura matemática necesitamos?}{}
  }
  \temaEsquema{Espacios vectoriales}{
    \conceptoEsquema{Definición y axiomas}{}
    \conceptoEsquema{Ejemplos: $\C^n$, matrices, polinomios}{}
    \conceptoEsquema{Subespacios vectoriales}{}
  }
  \temaEsquema{Dependencia lineal y bases}{
    \conceptoEsquema{Combinaciones lineales y generadores}{}
    \conceptoEsquema{Independencia lineal}{}
    \conceptoEsquema{Base y dimensión}{}
    \conceptoEsquema{Base canónica de $\C^n$ y base computacional $\{\ket{0},\ket{1}\}$}{}
  }
  \temaEsquema{Coordenadas y cambio de base}{
    \conceptoEsquema{Coordenadas de un vector en una base}{}
    \conceptoEsquema{Matriz de cambio de coordenadas}{}
  }
\end{esquemaExplorador}

% ------------------------------------------------------------
\tema{3}{Operadores lineales y representación matricial}
% ------------------------------------------------------------

\portada

\begin{info}
  \textbf{Motivación cuántica (anticipo de los Postulados II y V).}
  En mecánica cuántica, tanto la evolución temporal de un sistema (Postulado~V)
  como la acción de un observable (Postulado~II) se describen mediante
  \emph{transformaciones lineales}.
  Este tema construye la teoría de operadores lineales y su representación matricial,
  introduciendo además las clases de matrices (hermitianas, unitarias) que aparecen
  de forma central en los postulados cuánticos.
  Con estas herramientas será posible presentar la descomposición espectral y la SVD.
\end{info}

\begin{esquemaExplorador}
  \temaEsquema{Transformaciones lineales}{
    \conceptoEsquema{Definición de operador lineal}{}
    \conceptoEsquema{Núcleo e imagen}{}
    \conceptoEsquema{Isomorfismos y composición}{}
  }
  \temaEsquema{Representación matricial}{
    \conceptoEsquema{Matriz de un operador en una base}{}
    \conceptoEsquema{Cambio de base para operadores}{$A' = P^{-1}AP$}
    \conceptoEsquema{Traza y determinante}{}
  }
  \temaEsquema{Operaciones sobre matrices}{
    \conceptoEsquema{Transpuesta y conjugada}{}
    \conceptoEsquema{Adjunta (conjugada transpuesta)}{$A^\dagger = (\bar{A})^t$}
  }
  \temaEsquema{Clases especiales de matrices}{
    \conceptoEsquema{Matrices normales}{$AA^\dagger = A^\dagger A$}
    \conceptoEsquema{Matrices hermitianas}{$A^\dagger = A$}
    \conceptoEsquema{Matrices unitarias}{$UU^\dagger = U^\dagger U = I$}
    \conceptoEsquema{Propiedades y ejemplos cuánticos (Pauli, Hadamard)}{}
  }
  \temaEsquema{Valores y vectores propios}{
    \conceptoEsquema{Polinomio característico}{}
    \conceptoEsquema{Cálculo de vectores propios}{}
    \conceptoEsquema{Diagonalización}{$A = PDP^{-1}$}
  }
  \temaEsquema{Descomposición en valores singulares (SVD)}{
    \conceptoEsquema{Enunciado}{$A = U\Sigma V^\dagger$}
    \conceptoEsquema{Algoritmo de cálculo}{}
    \conceptoEsquema{Interpretación geométrica}{}
    \conceptoEsquema{Aplicación: análisis de entrelazamiento (anticipo)}{}
  }
\end{esquemaExplorador}

% ============================================================
%  BLOQUE 2: ESPACIOS CUÁNTICOS
% ============================================================
\bloque{2}{Espacios cuánticos}

% ------------------------------------------------------------
\tema{4}{Espacios de Hilbert}
% ------------------------------------------------------------

\portada

\begin{info}
  \textbf{Motivación cuántica (Postulados I, II y III).}
  \begin{itemize}
    \item \textbf{Postulado~I (completo):} El espacio de estados de un sistema cuántico
          es un \emph{espacio de Hilbert} $\mathcal{H}$, no un espacio vectorial genérico.
          La diferencia clave es el producto interno, que permite definir norma, ángulos y probabilidades.
    \item \textbf{Postulado~II:} Los observables físicos son operadores hermitianos,
          cuyas propiedades más profundas (valores propios reales, eigenvectores ortogonales)
          solo se comprenden con el producto interno.
    \item \textbf{Postulado~III:} La regla de Born, $p(\lambda_i)=|\braket{\lambda_i}{\psi}|^2$,
          es un producto interno. Sin esta estructura no hay probabilidades cuánticas.
  \end{itemize}
\end{info}

\begin{esquemaExplorador}
  \temaEsquema{Postulados I--III: enunciado preliminar}{
    \conceptoEsquema{Estado cuántico como vector unitario en $\mathcal{H}$}{}
    \conceptoEsquema{Observable como operador hermitiano}{}
    \conceptoEsquema{Regla de Born}{$p(\lambda_i)=|\langle\lambda_i|\psi\rangle|^2$}
  }
  \temaEsquema{Producto interno en $\C^n$}{
    \conceptoEsquema{Definición}{$\langle u,v\rangle = \sum_i \bar{u}_i v_i$}
    \conceptoEsquema{Propiedades: linealidad, simetría hermitiana, positividad}{}
    \conceptoEsquema{Norma inducida}{$\|v\|=\sqrt{\langle v,v\rangle}$}
  }
  \temaEsquema{Ortogonalidad y proyecciones}{
    \conceptoEsquema{Vectores ortogonales y complemento ortogonal}{}
    \conceptoEsquema{Proyección ortogonal sobre un subespacio}{}
    \conceptoEsquema{Proceso de Gram-Schmidt}{}
    \conceptoEsquema{Bases ortonormales}{}
  }
  \temaEsquema{Matrices especiales revisadas con producto interno}{
    \conceptoEsquema{Hermitianas: valores propios reales y eigenvectores ortogonales}{}
    \conceptoEsquema{Unitarias: preservan norma y producto interno}{}
    \conceptoEsquema{Teorema espectral para matrices normales}{}
  }
  \temaEsquema{Espacios de Hilbert (dimensión finita)}{
    \conceptoEsquema{Definición formal}{}
    \conceptoEsquema{Completitud e identidad de resolución}{}
    \conceptoEsquema{El espacio de un cúbit: $\C^2$}{}
  }
\end{esquemaExplorador}

% ------------------------------------------------------------
\tema{5}{Espacio dual y producto tensorial}
% ------------------------------------------------------------

\portada

\begin{info}
  \textbf{Motivación cuántica (Postulado IV).}
  El Postulado~IV establece que el espacio de estados de un sistema compuesto
  por dos subsistemas $A$ y $B$ es el \emph{producto tensorial}
  $\mathcal{H}_A \otimes \mathcal{H}_B$.
  Para entender esta construcción con rigor se necesita previamente el concepto
  de funcional lineal y espacio dual, que además es el fundamento matemático
  de la notación bra-ket que se desarrollará en el siguiente tema.
\end{info}

\begin{esquemaExplorador}
  \temaEsquema{Postulado IV: anticipo}{
    \conceptoEsquema{Sistema de dos cúbits}{$\mathcal{H}=\C^2\otimes\C^2\cong\C^4$}
    \conceptoEsquema{Base computacional}{$\{\ket{00},\ket{01},\ket{10},\ket{11}\}$}
    \conceptoEsquema{¿Qué estructura matemática describe la composición?}{}
  }
  \temaEsquema{Funcionales lineales y espacio dual}{
    \conceptoEsquema{Definición de funcional lineal $f:V\to\C$}{}
    \conceptoEsquema{Espacio dual $V^*$}{}
    \conceptoEsquema{Teorema de representación de Riesz-Fréchet}{}
    \conceptoEsquema{Base dual y base ortonormal}{}
  }
  \temaEsquema{Producto tensorial de espacios vectoriales}{
    \conceptoEsquema{Definición axiomática}{}
    \conceptoEsquema{Base del producto tensorial}{$\{e_i\otimes f_j\}$}
    \conceptoEsquema{Producto tensorial de matrices}{}
    \conceptoEsquema{Propiedades}{}
  }
  \temaEsquema{Sistemas cuánticos compuestos}{
    \conceptoEsquema{$n$ cúbits: espacio $\C^{2^n}$}{}
    \conceptoEsquema{Estados separables}{$\ket{\psi}\otimes\ket{\phi}$}
    \conceptoEsquema{Estados entrelazados: imposibilidad de factorizar}{}
    \conceptoEsquema{Estados de Bell}{}
  }
\end{esquemaExplorador}

% ============================================================
%  BLOQUE 3: OPERADORES Y DINÁMICAS CUÁNTICAS
% ============================================================
\bloque{3}{Operadores y dinámicas cuánticas}

% ------------------------------------------------------------
\tema{6}{Notación de Dirac y los postulados de la mecánica cuántica}
% ------------------------------------------------------------

\portada

\begin{info}
  \textbf{Motivación cuántica (los cinco postulados).}
  Hasta este punto hemos construido toda la maquinaria matemática necesaria:
  espacios de Hilbert, operadores hermitianos, producto interno y producto tensorial.
  La notación de Dirac es el \emph{lenguaje estándar} de la mecánica cuántica que unifica
  y simplifica todos estos elementos.
  En este tema se enuncia el formalismo completo mediante los \textbf{cinco postulados}
  de la mecánica cuántica, escritos en notación bra-ket.
\end{info}

\begin{esquemaExplorador}
  \temaEsquema{Notación bra-ket}{
    \conceptoEsquema{Ket: vector del espacio de Hilbert}{$\ket{\psi}\in\mathcal{H}$}
    \conceptoEsquema{Bra: elemento del espacio dual}{$\bra{\psi}\in\mathcal{H}^*$}
    \conceptoEsquema{Braket: producto interno}{$\braket{\phi}{\psi}\in\C$}
    \conceptoEsquema{Ketbra: operador externo}{$\ketbra{\psi}{\phi}$}
  }
  \temaEsquema{Operadores en notación Dirac}{
    \conceptoEsquema{Acción de un operador sobre un ket}{}
    \conceptoEsquema{Identidad de resolución}{$\sum_i\ketbra{i}{i}=I$}
    \conceptoEsquema{Descomposición espectral en Dirac}{$O=\sum_i\lambda_i\ketbra{\lambda_i}{\lambda_i}$}
  }
  \temaEsquema{Los cinco postulados}{
    \conceptoEsquema{P-I: Estado cuántico}{$\ket{\psi}\in\mathcal{H},\ \|\ket{\psi}\|=1$}
    \conceptoEsquema{P-II: Observable}{$O=O^\dagger$}
    \conceptoEsquema{P-III: Medición y regla de Born}{$p(\lambda_i)=|\braket{\lambda_i}{\psi}|^2$}
    \conceptoEsquema{P-IV: Sistema compuesto}{$\mathcal{H}=\mathcal{H}_A\otimes\mathcal{H}_B$}
    \conceptoEsquema{P-V: Evolución temporal}{$\ket{\psi(t)}=U\ket{\psi(0)}$}
  }
  \temaEsquema{Estados cuánticos en notación Dirac}{
    \conceptoEsquema{Cúbit}{$\alpha\ket{0}+\beta\ket{1},\ |\alpha|^2+|\beta|^2=1$}
    \conceptoEsquema{Medición en bases distintas}{}
    \conceptoEsquema{Estados de Bell y entrelazamiento máximo}{}
  }
\end{esquemaExplorador}

% ------------------------------------------------------------
\tema{7}{Operadores en computación cuántica}
% ------------------------------------------------------------

\portada

\begin{info}
  \textbf{Motivación cuántica (Postulado~V: evolución temporal).}
  El Postulado~V establece que la evolución de un sistema cuántico cerrado
  se describe mediante un operador unitario.
  En el modelo del circuito cuántico, las \emph{puertas cuánticas} son
  precisamente los operadores unitarios que implementan esa evolución de forma controlada.
  Este tema estudia las puertas cuánticas fundamentales como objetos del álgebra lineal
  y los representa geométricamente en la esfera de Bloch.
\end{info}

\begin{esquemaExplorador}
  \temaEsquema{Puertas de un cúbit}{
    \conceptoEsquema{Matrices de Pauli}{$\sigma_x,\sigma_y,\sigma_z$}
    \conceptoEsquema{Puerta Hadamard}{$H=\frac{1}{\sqrt{2}}\mqty(1&1\\1&-1)$}
    \conceptoEsquema{Puertas de fase}{$S,T$}
    \conceptoEsquema{Verificación de unitariedad}{}
  }
  \temaEsquema{Esfera de Bloch}{
  \conceptoEsquema{Parametrización de estados de un cúbit}{$\cos\frac{\theta}{2}\ket{0}+e^{i\phi}\sin\frac{\theta}{2}\ket{1}$}
  \conceptoEsquema{Puertas como rotaciones}{$R_x,R_y,R_z$}
  \conceptoEsquema{Interpretación geométrica de Hadamard y Pauli}{}
  }
  \temaEsquema{Puertas de múltiples cúbits}{
    \conceptoEsquema{Puerta CNOT}{$\ket{x,y}\mapsto\ket{x,x\oplus y}$}
    \conceptoEsquema{Operadores controlados en general}{}
    \conceptoEsquema{Producto tensorial de puertas}{}
  }
  \temaEsquema{Universalidad}{
    \conceptoEsquema{Concepto de conjunto universal de puertas}{}
    \conceptoEsquema{Ejemplo: $\{H,T,\mathrm{CNOT}\}$}{}
  }
\end{esquemaExplorador}

% ------------------------------------------------------------
\tema{8}{Información cuántica y sistemas abiertos}
% ------------------------------------------------------------

\portada

\begin{info}
  \textbf{Motivación cuántica (extensión del Postulado~III).}
  El formalismo de estado puro (Postulado~I) es insuficiente para describir
  un sistema cuántico que interactúa con su entorno o cuya preparación
  es estadísticamente imperfecta.
  La \emph{matriz de densidad} extiende el formalismo de los postulados para
  cubrir estos \textbf{estados mixtos}, que son los estados que
  aparecen en la práctica experimental y en los protocolos de información cuántica.
\end{info}

\begin{esquemaExplorador}
  \temaEsquema{Estados mixtos y operador densidad}{
    \conceptoEsquema{Estado puro vs. estado mixto}{}
    \conceptoEsquema{Definición del operador densidad}{$\rho=\sum_i p_i\ketbra{\psi_i}{\psi_i}$}
    \conceptoEsquema{Propiedades: $\rho^\dagger=\rho$, $\tr(\rho)=1$, $\rho\geq0$}{}
    \conceptoEsquema{Criterio de pureza}{$\tr(\rho^2)\leq1$}
  }
  \temaEsquema{Traza parcial y sistemas bipartitos}{
    \conceptoEsquema{Definición de traza parcial}{}
    \conceptoEsquema{Estado reducido de un subsistema}{}
    \conceptoEsquema{Entrelazamiento y matrices reducidas}{}
  }
  \temaEsquema{Canales cuánticos}{
    \conceptoEsquema{Operaciones cuánticas como mapas de matrices de densidad}{}
    \conceptoEsquema{Representación de Kraus}{$\mathcal{E}(\rho)=\sum_k K_k\rho K_k^\dagger$}
    \conceptoEsquema{Ejemplos: desfasamiento, amortiguamiento de amplitud}{}
  }
  \temaEsquema{Información cuántica}{
    \conceptoEsquema{Teorema de no-clonación}{}
    \conceptoEsquema{Entropía de von Neumann}{$S(\rho)=-\tr(\rho\log\rho)$}
    \conceptoEsquema{Fidelidad}{$F(\rho,\sigma)=\tr\!\sqrt{\sqrt{\rho}\,\sigma\sqrt{\rho}}$}
    \conceptoEsquema{Decoherencia y pérdida de coherencia}{}
  }
\end{esquemaExplorador}

\end{document}